% problem_id: S001-lemma-lowering-the-level
% source_id: S001 (Stacks Project)
% source_file: more-algebra.tex
% statement_locator: more-algebra.tex lines 35164--35180
% proof_locator: more-algebra.tex lines 35182--35367
% env: lemma
% id_note: label 在本来源内唯一
% pairing: tight (verified)
% status: complete (statement + author proof verbatim + reason + locators)
%
\begin{lemma}
\label{lemma-lowering-the-level}
Let $A \subset B \subset C$ be extensions of discrete valuation rings
with fractions fields $K \subset L \subset M$. Assume
\begin{enumerate}
\item $A$ has mixed characteristic $(0, p)$,
\item $A \subset B$ weakly unramified,
\item $B$ contains a primitive $p$th root of $1$,
\item $M/L$ is a degree $p$ extension of finite level $l > 0$,
\item $\kappa_A = \bigcap_{n \geq 1} \kappa_B^{p^n}$.
\end{enumerate}
Then there exists a finite separable extension $K_1$ of $K$
totally ramified with respect to $A$
such that either $K_1$ is a weak solution for $A \to C$, or the extension
$M_1/L_1$ is a degree $p$ extension of finite level
$\leq \max(0, l - 1, 2l - p)$.
\end{lemma}

\begin{proof}
Let $\pi \in A$ be a uniformizer.
Let $w \in B$ and $P \in B[t]$ be as in Lemma \ref{lemma-prepare} (for $B$).
Set $e_1 = \text{ord}_B(w)$, so that $w$ and $\pi^{e_1}$ are associates in $B$.
Pick $z \in M$ generating $M$ over $L$ with $\xi = P(z) \in K$
and $n$ such that $\pi^n\xi \in B$ as in the definition of the level
of $M$ over $L$, i.e., $l = n/e_1$.

\medskip\noindent
The proof of this lemma is completely similar to the proof of
Lemma \ref{lemma-characteristic-p-case}.
To explain what is going on, observe that
\begin{equation}
\label{equation-first-congruence}
P(z) \equiv z^p - z \bmod \pi^{-n + e_1}B
\end{equation}
for any $z \in L$ such that $\pi^{-n} P(z) \in B$ (use that $z$ has valuation
at worst $-n/p$ and the shape of the polynomial $P$). Moreover, we have
\begin{equation}
\label{equation-second-congruence}
\xi_1 + \xi_2 + w^p \xi_1 \xi_2 \equiv \xi_1 + \xi_2 \bmod \pi^{-2n + pe_1}B 
\end{equation}
for $\xi_1, \xi_2 \in \pi^{-n}B$. Finally, observe that
$n - e_1 = (l - 1)/e_1$ and $-2n + pe_1 = -(2l - p)e_1$.
Write $m = n - e_1 \max(0, l - 1, 2l - p)$. The above shows that doing
calculations in $\pi^{-n}B / \pi^{-n + m}B$ the polynomial $P$ behaves exactly
as the polynomial $z^p -  z$. This explains why the lemma is true
but we also give the details below.

\medskip\noindent
Assumption (4) implies that $\kappa_A$ is perfect. Observe that
$m \leq e_1$ and hence $A/\pi^m$ is annihilated by $w$ and hence $p$.
Thus we may choose compatible ring maps
$\overline{\sigma} : \kappa_A \to A/\pi^mA$ and
$\overline{\sigma} : \kappa_B \to B/\pi^mB$ as in
Lemma \ref{lemma-cohen}. We lift the second of these to a
map of sets $\sigma : \kappa_B \to B$. Then we can write
$$
\xi =
\sum\nolimits_{i = n, \ldots, n - m + 1} \sigma(\lambda_i) \pi^{-i} +
\pi^{-n + m)} b
$$
for some $\lambda_i \in \kappa_B$ and $b \in B$. Let
$$
I = \{i \in \{n, \ldots, n - m + 1\} \mid \lambda_i \in \kappa_A\}
$$
and
$$
J = \{j \in \{n, \ldots, n - m + 1\} \mid \lambda_i \not \in \kappa_A\}
$$
We will argue by induction on the size of the finite set $J$.

\medskip\noindent
The case $J = \emptyset$. Here for all $i \in \{n, \ldots, n - m + 1\}$
we have $\sigma(\lambda_i) = a_i + \pi^{n - m}b_i$ for some $a_i \in A$
and $b_i \in B$ by our choice of $\overline{\sigma}$. Thus
$\xi = \pi^{-n} a + \pi^{-n + m} b$ for some $a \in A$ and $b \in B$.
If $p | n$, then we write $a = a_0^p + \pi a_1$ for some $a_0, a_1 \in A$
(as the residue field of $A$ is perfect). Set $z_1 = - \pi^{-n/p} a_0$.
Note that $P(z_1) \in \pi^{-n}B$ and that $z + z_1 + w z z_1$ is an
element generating $M$ over $L$ (note that $wz_1 \not = -1$ as
$n < pe_1$). Moreover, by Lemma \ref{lemma-prepare} we have
$$
P(z + z_1 + w z z_1) = P(z) + P(z_1) + w^p P(z) P(z_1) \in K
$$
and by equations (\ref{equation-first-congruence}) and
(\ref{equation-second-congruence}) we have
$$
P(z) + P(z_1) + w^p P(z) P(z_1)
\equiv
\xi + z_1^p - z_1 \bmod \pi^{-n + m}B
$$
for some $b' \in B$. This contradict the minimality of $n$! Thus $p$
does not divide $n$. Consider the degree $p$ extension $K_1$ of $K$ given
by $P(y) = -\pi^{-n}a$. By Lemma \ref{lemma-extension-defined-by-nice-polynial}
this extension is separable and totally ramified with respect to $A$.
Thus $L_1 = L \otimes_K K_1$
is a field and $A_1 \subset B_1$ is weakly unramified
(Lemma \ref{lemma-weakly-unramified-goes-up-along-totally-ramified}).
By Lemma \ref{lemma-extension-defined-by-nice-polynial}
the ring $M_1 = M \otimes_K K_1$ is either a product of $p$ copies
of $L_1$ (in which case we are done) or a field extension of $L_1$
of degree $p$. Moreover, in the second case, either $C_1$ is weakly unramified
over $B_1$ (in which case we are done) or $M_1/L_1$ is degree $p$,
Galois, totally ramified with respect to $B_1$.
In this last case the extension $M_1/L_1$
is generated by the element $z + y + wzy$ and we see that
$P(z + y + wzy) \in L_1$ and
\begin{align*}
P(z + y + wzy)
& = P(z) + P(y) + w^p P(z) P(y) \\
& \equiv
\xi - \pi^{-n}a \bmod \pi^{-n + m}B_1 \\
& \equiv
0 \bmod \pi^{-n + m}B_1
\end{align*}
in exactly the same manner as above. By our choice of $m$ this
means exactly that $M_1/L_1$ has level at most $\max(0, l - 1, 2l - p)$.
From now on we assume that $J \not = \emptyset$.

\medskip\noindent
Suppose that $j', j \in J$ such that $j' = p^r j$ for some $r > 0$.
Then we set
$$
z_1 = - \sigma(\lambda_j) \pi^{-j} - \sigma(\lambda_j^p) \pi^{-pj} -
\ldots - \sigma(\lambda_j^{p^{r - 1}}) \pi^{-p^{r - 1}j}
$$
and we change $z$ into $z' = z + z_1 + wzz_1$. Observe that $z' \in M$
generates $M$ over $L$ and that we have
$\xi' = P(z') = P(z) + P(z_1) + wP(z)P(z_1) \in L$ with
$$
\xi' \equiv
\xi - \sigma(\lambda_j) \pi^{-j} + \sigma(\lambda_j^{p^r}) \pi^{-j'}
\bmod \pi^{-n + m}B
$$
by using equations (\ref{equation-first-congruence}) and
(\ref{equation-second-congruence}) as above. Writing
$$
\xi' = \sum\nolimits_{i = n, \ldots, n - m + 1} \sigma(\lambda'_i) \pi^{-i}
+ \pi^{-n + m}b'
$$
as before we find that
$\lambda'_i = \lambda_i$ for $i \not = j, j'$ and $\lambda'_j = 0$.
Thus the set $J$ has gotten smaller.
By induction on the size of $J$ we may assume there is no pair
$j, j'$ of $J$ such that $j'/j$ is a power of $p$.
(Please observe that in this procedure we may get thrown back into the case
that $J = \emptyset$ we treated above.)

\medskip\noindent
For $j \in J$ write $\lambda_j = \mu_j^{p^{r_j}}$ for some $r_j \geq 0$ and
$\mu_j \in \kappa_B$ which is not a $p$th power. This is possible by our
assumption (4). Let $j \in J$ be the unique index such that $j p^{-r_j}$
is maximal. (The index is unique by the result of the preceding paragraph.)
Choose $r > \max(r_j + 1)$ and such that $j p^{r - r_j} > n$ for $j \in J$.
Let $K_1/K$ be the extension of degree $p^r$, totally ramified
with respect to $A$, defined by $(\pi')^{p^r} = \pi$.
Observe that $\pi'$ is the uniformizer of the
corresponding discrete valuation ring $A_1 \subset K_1$.
Observe that $L_1 = L \otimes_K K_1$ is a field and $L_1/L$
is totally ramified with respect to $B$
(Lemma \ref{lemma-weakly-unramified-goes-up-along-totally-ramified}).
Computing in the integral closure $B_1$ we get
$$
\xi = \sum\nolimits_{i \in I} \sigma(\lambda_i) (\pi')^{-i p^r} +
\sum\nolimits_{j \in J} \sigma(\mu_j)^{p^{r_j}} (\pi')^{-j p^r} +
\pi^{-n + m} b_1
$$
for some $b_1 \in B_1$. Note that $\sigma(\lambda_i)$ for $i \in I$
is a $q$th power modulo $\pi^m$, i.e., modulo $(\pi')^{m p^r}$.
Hence we can rewrite the above as
$$
\xi = \sum\nolimits_{i \in I} x_i^{p^r} (\pi')^{-i p^r} +
\sum\nolimits_{j \in J} \sigma(\mu_j)^{p^{r_j}} (\pi')^{- j p^r}
+ \pi^{-n + m}b_1
$$
Similar to our choice in the previous paragraph we set
\begin{align*}
z_1 & - \sum\nolimits_{i \in I}
\left(x_i (\pi')^{-i} + \ldots + x_i^{p^{r - 1}} (\pi')^{-i p^{r - 1}}\right)
\\
& - \sum\nolimits_{j \in J}
\left(
\sigma(\mu_j) (\pi')^{- j p^{r - r_j}}
+ \ldots +
\sigma(\mu_j)^{p^{r_j - 1}} (\pi')^{- j p^{r - 1}}
\right)
\end{align*}
and we change our choice of $z$ into $z' = z + z_1 + wzz_1$.
Then $z'$ generates $M_1$ over $L_1$ and
$\xi' = P(z') = P(z) + P(z_1) + w^p P(z) P(z_1) \in L_1$
and a calculation shows that
$$
\xi' \equiv
\sum\nolimits_{i \in I} x_i (\pi')^{-i} +
\sum\nolimits_{j \in J} \sigma(\mu_j) (\pi')^{- j p^{r - r_j}} +
(\pi')^{(-n + m)p^r}b'_1
$$
for some $b'_1 \in B_1$. There is a unique $j$ such that $j p^{r - r_j}$
is maximal and $j p^{r - r_j}$ is bigger than $i \in I$. If
$j p^{r - r_j} \leq (n - m)p^r$ then the level of the extension $M_1/L_1$
is less than $\max(0, l - 1, 2l - p)$. If not, then, as $p$ divides
$j p^{r - r_j}$, we see that $M_1 / L_1$ falls into case (C) of
Lemma \ref{lemma-extension-defined-by-nice-polynial}.
This finishes the proof.
\end{proof}
